{\scriptsize
        \begin{tabularx}{\textwidth}{|p{0.3\columnwidth}|X|}
            \hline 
            \textbf{Notions et contenus} & \textbf{Capacités exigibles} \\
            \hline 
            \multicolumn{2}{|l|}{\textbf{Approche ondulatoire de la mécanique quantique}} \\
            \hline
            \multicolumn{2}{|l|}{\textbf{6.5.1. Amplitude de probabilité}} \\
            \hline
            Fonction d'onde $\psi(x, t)$ associée à une particule dans un problème unidimensionnel. Densité linéique de probabilité de présence.  & 
            Normaliser une fonction d'onde. Relier qualitativement la fonction d'onde à la notion d'orbitale en chimie. \\
            \hline 
            Normaliser une fonction d'onde. Relier qualitativement la fonction d'onde à la notion d'orbitale en chimie.   & 
            Relier la superposition de fonctions d'ondes à la description d'une expérience d'interférences entre particules. \\
        \hline 
        \multicolumn{2}{|l|}{\textbf{6.5.2 Équation de Schrödinger pour une particule libre}} \\
        \hline    
        Équation de Schrödinger. & 
        Utiliser l'équation de Schrödinger fournie. \\
        \hline
         États stationnaires. &
        Associer les états stationnaires aux états d'énergie déterminée. Établir et utiliser la forme $\psi(\mathrm{x}, \mathrm{t})=\phi(\mathrm{x}) \exp (-\mathrm{iEt} / \hbar)$ pour la fonction d'onde d'un état stationnaire et l'associer à la relation de Planck-Einstein.
Distinguer l'onde associée à un état stationnaire en mécanique quantique d'une onde stationnaire au sens usuel de la physique des ondes.\\
        \hline
        Paquet d'ondes associé à une particule libre. Relation $\Delta k_x \Delta x \geq 1 / 2$. & Paquet d'ondes associé à une particule libre. Relation $\Delta k_x \Delta x \geq 1 / 2$. \\
        \hline
        Courant de probabilité associé à une particule libre. &
        Courant de probabilité associé à une particule libre. \\
        \hline
        \end{tabularx}
\vspace{0.5cm}
}
